%% Proposed solution: How do you plan to go about solving this problem? It is okay to not know
%% how the machine learning models work at this point in the class, but you should be starting
%% to get some idea based on the lectures and assigned readings. What kind of features and
%% target labels do you have? What kind of models might you try? You may not propose simple
%% linear regression models for this project. Are there any existing solutions to your problem?
%% You must indicate 2-5 sources (research papers, books, machine learning challenges online)
%% from where your solution was inspired. Put some thought into how you would approach this.
%% How will you know if the model is good? How will you evaluate it? Also, share the libraries
%% you intend to use for the project

\section{Solution}
U-Net is a convolutional neural network (CNN) architecture designed for image segmentation tasks. This architecture consists of an encoder and decoder. The encoder repeated applies convolutions and downsamples the image. It learns increasingly abstract features.  The decoder upsamples these features back up to the original resolution of the input. Its job is to figure out the exact location of features. In the context of our project, encoded early layers will detect edges and features, while deeper layers can learn patterns associated with tumors. \cite{unet}


Each MRI scan consists of four images: T1, contrast-enhanced T1 (T1CE), T2, and T2-FLAIR. These images will be stacked as separate input channels, allowing the model to use information from each image concurrently in its prediction. 

The target labels will consist of the already labeled segmentation masks provided with the dataset. Therefore, as the labels are pixel level, the model will also perform pixel-wise classification, assigning each pixel in the U-Net slice to either the background or one of the annotated tumor regions. As such, the model will be trained as a 2D multiclass U-Net on axial slices.

Each MRI modality will be normalized separately since MRI intensities do not have a universal scale. 2D axial slices will be extracted, and T1/T1CE/T2/FLAIR will be stacked as four channels. Augmentation will then be applied during training. 

The split will be done by patient/subject, not by individual slice. Otherwise adjacent slices from the same patient could appear in both training and testing which will artificially inflate results.

Dice score will be used as the main segmentation metric. It will be evaluated for individual tumor regions. We will compare cross-entropy loss versus Dice-based loss or a combined Dice + cross-entropy loss.

Relevant libraries are PyTorch, NumPy, NiBabel (for reading nii.gz data files).
